离散谱是空间谱方程在 $\mathbb C\mathbin{\backslash}\Sigma$ 中存在非零 $L^2(\mathbb R)$ 解的谱点集合. 设三场满足命题 \ref{prop:evolution} 的条件, 并假设透射系数在各自解析域中只有有限个零点, 且这些零点均为简单零点.

\begin{proposition}[四点离散谱轨道]\label{prop:quartet}
离散谱恰为 $s_{22}$ 在 $D_+$ 中及 $s_{11}$ 在 $D_-$ 中的零点. 若 $\zeta\in D_+$ 是 $s_{22}$ 的简单零点, 则
\begin{equation}\label{eq:quartet}
\left\{\zeta,\conj{\zeta},-\dfrac{q_0^2}{\zeta},-\dfrac{q_0^2}{\conj{\zeta}}\right\}
\end{equation}
构成四点互异的完整轨道. 其中 $\zeta$ 与 $h(\conj{\zeta})$ 属于 $D_+$, 其余两点属于 $D_-$. 每个轨道选取外部上半平面的代表 $\zeta_j$, $j=1,\ldots,N$, 再令
\begin{equation}\label{eq:discrete-indexing}
\zeta_{N+j}=h(\conj{\zeta_j}),\qquad
Z=\{\zeta_j,\conj{\zeta_j}:j=1,\ldots,2N\}.
\end{equation}
每个 $D_+$ 零点满足非零列相关关系
\begin{equation}\label{eq:norming-definition}
\phi_{+,2}(\zeta_j)=b_j\phi_{-,1}(\zeta_j),\qquad
c_j=\dfrac{b_j}{s_{22}'(\zeta_j)},\qquad
C_j=c_je^{2\ii\theta(\zeta_j)}.
\end{equation}
这里 $j=1,\ldots,2N$, $b_j,c_j$ 与 $x,t$ 无关, 撇号表示对单值化变量 $z$ 求导.
\end{proposition}
\begin{proof}
在 $D_+$ 中, $\operatorname{Im}\lambda>0$, 因而
\[
\phi_{-,1}\sim Y_1e^{\ii\theta}\quad(x\to-\infty),\qquad
\phi_{+,2}\sim Y_2e^{-\ii\theta}\quad(x\to+\infty)
\]
分别指数衰减. 任意另一个半轴 $L^2$ 解与该衰减列的 Wronskian 为常数; Cauchy--Schwarz 不等式又保证该 Wronskian 在半轴上可积, 因而此常数只能为零. 两端的 $L^2$ 解空间因此各为一维, 一个解在两端均属 $L^2$ 当且仅当这两列线性相关. 由 \eqref{eq:wronskian}, 这又等价于 $s_{22}(\zeta)=0$. 在 $D_-$ 中交换两列即得 $s_{11}$ 的结论.

共轭与换叶各自交换两个解析域, 二者复合保留解析域. 由 \eqref{eq:sconj}--\eqref{eq:ssheet},
\[
s_{22}(\zeta)=0
\quad\Longrightarrow\quad
s_{11}(\conj{\zeta})=s_{11}(h(\zeta))=s_{22}(h(\conj{\zeta}))=0.
\]
$h'(z)=\dfrac{q_0^2}{z^2}\ne0$ 保持零点的简单性. 共轭的不动点位于实轴, 换叶的不动点为 $\pm\ii q_0$, 而 $h(\conj{z})=z$ 将导致 $|z|^2=-q_0^2$. 因而 $\Sigma$ 外的四点互异, 且每个轨道恰有一个外部上半平面的代表.

Jost 列非零, 故列相关常数 $b_j\ne0$. 对 $\phi_{+,2}=b_j\phi_{-,1}$ 分别求 $x,t$ 导数并使用共同的 Lax 对, 得到
\[
(b_j)_x\phi_{-,1}=0,\qquad (b_j)_t\phi_{-,1}=0.
\]
于是 $b_j$ 与 $x,t$ 无关. 透射系数及其导数也不随 $x,t$ 改变, 从而得到 \eqref{eq:norming-definition} 中的常数性质.
\end{proof}

\begin{proposition}[简单零点的留数及常数约化]\label{prop:respair}
在式 \eqref{eq:norming-definition} 的条件下, 相应解析商满足
\begin{align}
\Res_{z=\zeta_j}\dfrac{\mu_{+,2}(z)}{s_{22}(z)}
&=C_j\mu_{-,1}(\zeta_j),\label{eq:resplus}\\
\Res_{z=\conj{\zeta_j}}\dfrac{\mu_{+,1}(z)}{s_{11}(z)}
&=-\conj{C_j}\mu_{-,2}(\conj{\zeta_j}).\label{eq:resminus}
\end{align}
若 $\widehat\zeta=h(\conj{\zeta})$ 为同一轨道的另一个 $D_+$ 点, 则
\begin{equation}\label{eq:cpair}
b_{\widehat\zeta}=-\conj{b_\zeta},\qquad
c_{\widehat\zeta}=-\dfrac{q_0^2}{(\conj{\zeta})^2}\conj{c_\zeta}.
\end{equation}
相应时空留数系数满足
\begin{equation}\label{eq:Cpair}
C_{\widehat\zeta}=-\dfrac{q_0^2}{(\conj{\zeta})^2}\conj{C_\zeta}.
\end{equation}
\end{proposition}
\begin{proof}
由 $\phi_{+,2}=b_j\phi_{-,1}$ 及 $\phi_1=\mu_1e^{\ii\theta}$、$\phi_2=\mu_2e^{-\ii\theta}$, 得到
\[
\mu_{+,2}(\zeta_j)=b_je^{2\ii\theta(\zeta_j)}\mu_{-,1}(\zeta_j).
\]
在简单零点处展开,
\[
s_{22}(z)=s_{22}'(\zeta_j)(z-\zeta_j)+O((z-\zeta_j)^2),\qquad
s_{22}'(\zeta_j)\ne0.
\]
因此
\[
\Res_{z=\zeta_j}\dfrac{\mu_{+,2}(z)}{s_{22}(z)}
=\lim\limits_{z\to\zeta_j}\dfrac{(z-\zeta_j)\mu_{+,2}(z)}{s_{22}(z)}
=\dfrac{\mu_{+,2}(\zeta_j)}{s_{22}'(\zeta_j)},
\]
代入列相关关系即得 \eqref{eq:resplus}. 对列相关关系应用 \eqref{eq:jost-conjugate-columns}, 在 $\conj{\zeta_j}$ 处得到
\[
\phi_{+,1}=-\conj{b_j}\phi_{-,2},\qquad
s_{11}'(\conj{\zeta_j})=\conj{s_{22}'(\zeta_j)}.
\]
这里的导数关系也可由上述 Taylor 展开取共轭直接得到. 又因 $\theta(\conj{\zeta_j})=\conj{\theta(\zeta_j)}$,
\[
\dfrac{\mu_{+,1}(\conj{\zeta_j})}{s_{11}'(\conj{\zeta_j})}
=-\conj{c_j}e^{-2\ii\theta(\conj{\zeta_j})}
\mu_{-,2}(\conj{\zeta_j})
=-\conj{C_j}\mu_{-,2}(\conj{\zeta_j}),
\]
即为 \eqref{eq:resminus}.

再在 $z=\conj{\zeta}$ 使用换叶关系. 它给出
\[
\phi_{+,2}(\widehat\zeta)
=\dfrac{\ii\conj{\zeta}}{q_0}\phi_{+,1}(\conj{\zeta})
=-\conj{b_\zeta}\dfrac{\ii\conj{\zeta}}{q_0}\phi_{-,2}(\conj{\zeta})
=-\conj{b_\zeta}\phi_{-,1}(\widehat\zeta).
\]
对恒等式 $s_{22}(h(z))=s_{11}(z)$ 求导, 有
\[
s_{22}'(\widehat\zeta)h'(\conj{\zeta})=\conj{s_{22}'(\zeta)},\qquad
h'(\conj{\zeta})=\dfrac{q_0^2}{(\conj{\zeta})^2}.
\]
链式因子也可从两次展开直接看出:
\[
\begin{aligned}
h(z)-\widehat\zeta
&=h'(\conj{\zeta})(z-\conj{\zeta})+O((z-\conj{\zeta})^2),\\
s_{22}(h(z))
&=s_{22}'(\widehat\zeta)h'(\conj{\zeta})(z-\conj{\zeta})+O((z-\conj{\zeta})^2).
\end{aligned}
\]
因此
\[
c_{\widehat\zeta}
=\dfrac{-\conj{b_\zeta}}{s_{22}'(\widehat\zeta)}
=-h'(\conj{\zeta})\conj{c_\zeta},
\]
即为 \eqref{eq:cpair}. 最后 $\theta(\widehat\zeta)=-\conj{\theta(\zeta)}$, 故
\[
C_{\widehat\zeta}
=-\dfrac{q_0^2}{(\conj{\zeta})^2}\conj{c_\zeta}
e^{-2\ii\conj{\theta(\zeta)}}
=-\dfrac{q_0^2}{(\conj{\zeta})^2}\conj{C_\zeta},
\]
得到 \eqref{eq:Cpair}. 对全部 $D_+$ 零点应用 \eqref{eq:resplus}--\eqref{eq:resminus}, 即给出完整四点轨道的留数条件.
\end{proof}


\subsubsection*{本部分全部谱点与留数公式的逐式计算}
\paragraph{式 \eqref{eq:quartet} 之前: 两列的指数衰减.}
固定 $t$ 与谱点, 原文相位直接给出
\begin{align*}
e^{\ii\theta}&=e^{-\ii\lambda x}e^{\ii\lambda Ft},&
|e^{\ii\theta}|&=e^{x\operatorname{Im}\lambda}
 e^{-t\operatorname{Im}(\lambda F)},\\
e^{-\ii\theta}&=e^{\ii\lambda x}e^{-\ii\lambda Ft},&
|e^{-\ii\theta}|&=e^{-x\operatorname{Im}\lambda}
 e^{t\operatorname{Im}(\lambda F)}.
\end{align*}
在 $D_+$ 中 $\operatorname{Im}\lambda>0$.
第一行随 $x\to-\infty$ 衰减, 第二行随 $x\to+\infty$ 衰减.
对应的 $\mu$ 列在归一化半轴有界, 因而
\begin{align*}
\int_{-\infty}^0\|\phi_{-,1}(x)\|_1^2\,\dd x
&\le\|\mu_{-,1}\|_{L^\infty(-\infty,0)}^2
 e^{-2t\operatorname{Im}(\lambda F)}
 \int_{-\infty}^0 e^{2x\operatorname{Im}\lambda}\,\dd x\\
&=\dfrac{\|\mu_{-,1}\|_{L^\infty(-\infty,0)}^2
 e^{-2t\operatorname{Im}(\lambda F)}}{2\operatorname{Im}\lambda}<\infty,\\
\int_0^{+\infty}\|\phi_{+,2}(x)\|_1^2\,\dd x
&\le\|\mu_{+,2}\|_{L^\infty(0,+\infty)}^2
 e^{2t\operatorname{Im}(\lambda F)}
 \int_0^{+\infty} e^{-2x\operatorname{Im}\lambda}\,\dd x\\
&=\dfrac{\|\mu_{+,2}\|_{L^\infty(0,+\infty)}^2
 e^{2t\operatorname{Im}(\lambda F)}}{2\operatorname{Im}\lambda}<\infty.
\end{align*}
这里把半轴在 $0$ 分开只是选择积分区间, 没有对变量换元.
两个指数积分分别由
\[
\left.\dfrac{e^{2x\operatorname{Im}\lambda}}{2\operatorname{Im}\lambda}
\right|_{-\infty}^0
=\dfrac{1}{2\operatorname{Im}\lambda},
\qquad
\left.-\dfrac{e^{-2x\operatorname{Im}\lambda}}{2\operatorname{Im}\lambda}
\right|_0^{+\infty}
=\dfrac{1}{2\operatorname{Im}\lambda}
\]
计算得到.
在 $D_-$ 中 $\operatorname{Im}\lambda<0$,
$\phi_{+,1}$ 在正端平方可积, $\phi_{-,2}$ 在负端平方可积.
把上面的积分指数符号与端点同时反转,
两个积分的值都变为 $-\dfrac{1}{2\operatorname{Im}\lambda}>0$.

\paragraph{平方可积条件怎样等价于透射零点.}
第4部分对行列式的求导计算适用于同一系统的任意两列,
因此任意两个空间解的 Wronskian 都是常数.
如果一个是衰减 Jost 列, 另一个也是同一半轴上的 $L^2$ 解,
二阶行列式绝对值不超过两列欧氏范数的乘积.
将该不等式在半轴上积分, 再用积分型 Cauchy--Schwarz 不等式,
得到的上界是两列各自平方积分的平方根之积, 因而有限.
这里没有对另一解另起辅助记号; 所用不等式就是行列式与范数的基本估计.
一个非零常数在无限长半轴上的绝对值积分为无穷,
所以这个 Wronskian 常数只能为零.
于是任意半轴 $L^2$ 解都与对应 Jost 衰减列线性相关.
两列在一处相关后, 空间线性初值问题的唯一性使同一个关系在全轴成立.

因此 $D_+$ 内存在全轴 $L^2$ 解, 当且仅当
\[
\det(\phi_{-,1},\phi_{+,2})=0.
\]
用原文 Wronskian 公式直接替换,
\[
\det(\phi_{-,1},\phi_{+,2})=\gamma s_{22},
\qquad \gamma\ne0\quad\hbox{在 }D_+,
\]
可知这与 $s_{22}=0$ 等价.
反方向也具体成立: 如果 $s_{22}=0$, 两个非零列相关,
$\phi_{-,1}$ 在负端衰减, 在正端又是 $\phi_{+,2}$ 的非零常数倍,
所以它在全轴平方可积.
在 $D_-$ 中用
\[
\det(\phi_{+,1},\phi_{-,2})=\gamma s_{11}
\]
得到同样结论.
没有把一般非实谱上的指数增长列误认为 $L^2$ 列.

\paragraph{式 \eqref{eq:quartet}: 四点的零点类型与区域.}
由共轭约化先有
\[
s_{11}(\conj{\zeta})=\conj{s_{22}(\zeta)}=0.
\]
再由换叶约化,
\[
s_{11}(h(\zeta))=s_{22}(\zeta)=0,
\qquad
s_{22}(h(\conj{\zeta}))=s_{11}(\conj{\zeta})=0.
\]
两种映射分别交换 $D_+$ 与 $D_-$,
所以四点的归属逐个是
\[
\zeta\in D_+,
\qquad \conj{\zeta}\in D_-,
\qquad h(\zeta)\in D_-,
\qquad h(\conj{\zeta})\in D_+.
\]
共轭两次与换叶两次都回到原点, 且两种变换对易,
所以再对这四点作任意一次约化不会产生第五个点.

可能的重合只有以下三类, 其余情形可再作一次约化化到这三类:
\begin{align*}
\zeta=\conj{\zeta}
&\ \Longrightarrow\ \operatorname{Im}\zeta=0,\\
\zeta=h(\zeta)
&\ \Longrightarrow\ \zeta=-\dfrac{q_0^2}{\zeta}
\ \Longrightarrow\ \zeta^2=-q_0^2
\ \Longrightarrow\ \zeta=\pm\ii q_0,\\
\zeta=h(\conj{\zeta})
&\ \Longrightarrow\ \zeta\conj{\zeta}=-q_0^2
\ \Longrightarrow\ |\zeta|^2=-q_0^2.
\end{align*}
第一类位于实连续谱上, 第二类位于分支点上,
第三类因左侧非负而右侧为负不可能成立.
所以连续谱外四点互异.

\paragraph{式 \eqref{eq:discrete-indexing}: 为什么每轨道只有一个所选代表.}
$D_+$ 的两块区域是外部上半平面与内部下半平面.
若 $\zeta$ 已位于外部上半平面, 它就是代表.
若 $\zeta$ 位于内部下半平面, 则
\[
|h(\conj{\zeta})|=\dfrac{q_0^2}{|\zeta|}>q_0,
\qquad
\operatorname{Im}h(\conj{\zeta})
=-\dfrac{q_0^2}{|\zeta|^2}\operatorname{Im}\zeta>0.
\]
所以同一轨道的另一个 $D_+$ 点恰好是外部上半平面代表.
另外两点在 $D_-$, 因而不可能再成为这种代表.
枚举 $N$ 个轨道后, 前 $N$ 个点是这些代表,
后 $N$ 个点正好由 $h(\conj{\zeta_j})$ 补齐,
再加上共轭点就得到 $4N$ 个完整离散谱点.

\paragraph{式 \eqref{eq:norming-definition}: 非零常数及其时空恒定性.}
Jost 列不能在任何有限点等于零向量.
否则该点的零初值会使线性空间方程的解恒为零,
与归一化端的非零 $Y_j$ 矛盾.
两个非零列行列式为零时, 存在唯一的非零相关倍数,
因此 $b_j$ 有定义且非零.
逐项求空间导数,
\begin{align*}
(\phi_{+,2})_x
&=(b_j)_x\phi_{-,1}+b_j(\phi_{-,1})_x\\
&=(b_j)_x\phi_{-,1}+b_jU\phi_{-,1},\\
U\phi_{+,2}&=b_jU\phi_{-,1}.
\end{align*}
相减后得 $(b_j)_x\phi_{-,1}=0$,
而非零列保证 $(b_j)_x=0$.
时间方向同样有
\begin{align*}
(\phi_{+,2})_t
&=(b_j)_t\phi_{-,1}+b_jV\phi_{-,1},\\
V\phi_{+,2}&=b_jV\phi_{-,1},
\end{align*}
从而 $(b_j)_t=0$.
由于透射零点本身已由演化命题固定,
这里不出现零点随时间移动的链式项.
$s_{22}'(\zeta_j)$ 也恒定, 所以 $c_j$ 恒定.
而 $C_j$ 包含指数, 它实际满足
\[
(C_j)_x=-2\ii\lambda(\zeta_j)C_j,
\qquad
(C_j)_t=2\ii\lambda(\zeta_j)F(k(\zeta_j))C_j.
\]
这是直接对 $c_je^{2\ii\theta(\zeta_j)}$ 求导得到的,
与原文只称 $b_j,c_j$ 为常数完全一致.

\paragraph{式 \eqref{eq:resplus}: 指数转换与分子分母的 Taylor 计算.}
将原 Jost 列与修正列的定义逐个代入,
\[
\mu_{+,2}(\zeta_j)e^{-\ii\theta(\zeta_j)}
=b_j\mu_{-,1}(\zeta_j)e^{\ii\theta(\zeta_j)}.
\]
两侧乘 $e^{\ii\theta(\zeta_j)}$ 后,
\[
\mu_{+,2}(\zeta_j)
=b_je^{2\ii\theta(\zeta_j)}\mu_{-,1}(\zeta_j).
\]
现在同时展开分子与分母:
\begin{align*}
\mu_{+,2}(z)
&=\mu_{+,2}(\zeta_j)+(z-\zeta_j)\mu_{+,2}'(\zeta_j)
+O((z-\zeta_j)^2),\\
s_{22}(z)
&=(z-\zeta_j)\left(s_{22}'(\zeta_j)
+\dfrac{s_{22}''(\zeta_j)}{2}(z-\zeta_j)
+O((z-\zeta_j)^2)\right).
\end{align*}
由于 $s_{22}'(\zeta_j)\ne0$, 对括号求倒数并相乘, 得
\begin{align*}
\dfrac{1}{s_{22}(z)}
&=\dfrac{1}{s_{22}'(\zeta_j)(z-\zeta_j)}
-\dfrac{s_{22}''(\zeta_j)}{2s_{22}'(\zeta_j)^2}
+O(z-\zeta_j),\\
\dfrac{\mu_{+,2}(z)}{s_{22}(z)}
&=\dfrac{\mu_{+,2}(\zeta_j)}{s_{22}'(\zeta_j)(z-\zeta_j)}
+\dfrac{\mu_{+,2}'(\zeta_j)}{s_{22}'(\zeta_j)}\\
&\quad-\dfrac{\mu_{+,2}(\zeta_j)s_{22}''(\zeta_j)}{2s_{22}'(\zeta_j)^2}
+O(z-\zeta_j).
\end{align*}
倒数展开的一次系数也可以通过乘回验证:
常数项为 1, 一次项是
$\dfrac{s_{22}''}{2s_{22}'}-\dfrac{s_{22}''}{2s_{22}'}=0$.
因此极点系数正是
\[
\dfrac{\mu_{+,2}(\zeta_j)}{s_{22}'(\zeta_j)}
=\dfrac{b_j}{s_{22}'(\zeta_j)}e^{2\ii\theta(\zeta_j)}
\mu_{-,1}(\zeta_j)=C_j\mu_{-,1}(\zeta_j).
\]
分子在该点非零, 所以这里不是可去奇点.

\paragraph{式 \eqref{eq:resminus}: 共轭产生的负号逐列检查.}
原相关式取共轭后为
\[
\conj{\phi_{+,2}(\zeta_j)}
=\conj{b_j}\conj{\phi_{-,1}(\zeta_j)}.
\]
左乘 $\ii\sigma_2$,
\begin{align*}
\phi_{+,1}(\conj{\zeta_j})
&=\conj{b_j}\ii\sigma_2\conj{\phi_{-,1}(\zeta_j)},\\
\phi_{-,2}(\conj{\zeta_j})
&=-\ii\sigma_2\conj{\phi_{-,1}(\zeta_j)}.
\end{align*}
第二行的 $-\ii$ 与第一行的 $+\ii$ 相反,
因此
\[
\phi_{+,1}(\conj{\zeta_j})
=-\conj{b_j}\phi_{-,2}(\conj{\zeta_j}).
\]
改写成修正列,
\[
\mu_{+,1}(\conj{\zeta_j})
=-\conj{b_j}e^{-2\ii\theta(\conj{\zeta_j})}
\mu_{-,2}(\conj{\zeta_j}).
\]
现在被除的列是第一列, 所以指数是负两倍.

\paragraph{共轭零点处的导数由 Taylor 系数比较得到.}
不能把非解析映射 $z\mapsto\conj{z}$ 当作解析函数套用复导数链式法则.
正确计算从
$s_{11}(z)=\conj{s_{22}(\conj{z})}$ 的局部展开开始:
\begin{align*}
s_{22}(\conj{z})
&=s_{22}'(\zeta_j)(\conj{z}-\zeta_j)
+O(|\conj{z}-\zeta_j|^2),\\
s_{11}(z)
&=\conj{s_{22}'(\zeta_j)}(z-\conj{\zeta_j})
+O(|z-\conj{\zeta_j}|^2).
\end{align*}
因此
\[
s_{11}'(\conj{\zeta_j})=\conj{s_{22}'(\zeta_j)}\ne0.
\]
又有
\begin{align*}
\conj{C_j}
&=\conj{c_j}\conj{e^{2\ii\theta(\zeta_j)}}\\
&=\conj{c_j}e^{-2\ii\conj{\theta(\zeta_j)}}
=\conj{c_j}e^{-2\ii\theta(\conj{\zeta_j})}.
\end{align*}
所以第二个简单极点的留数为
\begin{align*}
\dfrac{\mu_{+,1}(\conj{\zeta_j})}{s_{11}'(\conj{\zeta_j})}
&=-\dfrac{\conj{b_j}}{\conj{s_{22}'(\zeta_j)}}
 e^{-2\ii\theta(\conj{\zeta_j})}\mu_{-,2}(\conj{\zeta_j})\\
&=-\conj{C_j}\mu_{-,2}(\conj{\zeta_j}).
\end{align*}
共轭后的简单性也在非零的一次 Taylor 系数中同时得到证明.

\paragraph{式 \eqref{eq:cpair}: 先算列相关常数, 再算导数.}
在 $\conj{\zeta}$ 处使用两条换叶列关系,
\begin{align*}
\phi_{+,2}(\widehat\zeta)
&=\dfrac{\ii\conj{\zeta}}{q_0}\phi_{+,1}(\conj{\zeta}),\\
\phi_{-,1}(\widehat\zeta)
&=\dfrac{\ii\conj{\zeta}}{q_0}\phi_{-,2}(\conj{\zeta}).
\end{align*}
把共轭点的列相关关系代入第一行,
\begin{align*}
\phi_{+,2}(\widehat\zeta)
&=-\conj{b_\zeta}\dfrac{\ii\conj{\zeta}}{q_0}
\phi_{-,2}(\conj{\zeta})\\
&=-\conj{b_\zeta}\phi_{-,1}(\widehat\zeta).
\end{align*}
右侧列非零, 所以相关倍数唯一,
$b_{\widehat\zeta}=-\conj{b_\zeta}$.

换叶的差值可以精确通分:
\begin{align*}
h(z)-h(\conj{\zeta})
&=-\dfrac{q_0^2}{z}+\dfrac{q_0^2}{\conj{\zeta}}
=\dfrac{q_0^2(z-\conj{\zeta})}{z\conj{\zeta}},\\
\lim\limits_{z\to\conj{\zeta}}
\dfrac{h(z)-h(\conj{\zeta})}{z-\conj{\zeta}}
&=\dfrac{q_0^2}{(\conj{\zeta})^2}.
\end{align*}
对 $s_{22}(h(z))=s_{11}(z)$ 写差商,
\[
\dfrac{s_{22}(h(z))-s_{22}(\widehat\zeta)}{h(z)-\widehat\zeta}
\dfrac{h(z)-\widehat\zeta}{z-\conj{\zeta}}
=\dfrac{s_{11}(z)-s_{11}(\conj{\zeta})}{z-\conj{\zeta}}.
\]
取极限后,
\[
s_{22}'(\widehat\zeta)\dfrac{q_0^2}{(\conj{\zeta})^2}
=\conj{s_{22}'(\zeta)},
\qquad
s_{22}'(\widehat\zeta)
=\dfrac{(\conj{\zeta})^2}{q_0^2}\conj{s_{22}'(\zeta)}.
\]
两因子都不为零, 所以另一 $D_+$ 零点仍然简单.
把结果代入 $c$ 的定义,
\begin{align*}
c_{\widehat\zeta}
&=\dfrac{-\conj{b_\zeta}}
{\dfrac{(\conj{\zeta})^2}{q_0^2}\conj{s_{22}'(\zeta)}}\\
&=-\dfrac{q_0^2}{(\conj{\zeta})^2}
\conj{\left(\dfrac{b_\zeta}{s_{22}'(\zeta)}\right)}
=-\dfrac{q_0^2}{(\conj{\zeta})^2}\conj{c_\zeta}.
\end{align*}
这里分母导数乘逆因子, 但再除以导数使留数系数乘正向因子.
所有撇号都对原文的 $z$ 求导, 没有改用 $k$ 导数.

\paragraph{式 \eqref{eq:Cpair}: 相位与四个留数的完整检查.}
由 $k$ 在换叶下不变而 $\lambda$ 变号,
\[
\theta(\widehat\zeta)=-\conj{\theta(\zeta)},
\qquad
e^{2\ii\theta(\widehat\zeta)}
=e^{-2\ii\conj{\theta(\zeta)}}
=\conj{e^{2\ii\theta(\zeta)}}.
\]
所以
\begin{align*}
C_{\widehat\zeta}
&=c_{\widehat\zeta}e^{2\ii\theta(\widehat\zeta)}\\
&=-\dfrac{q_0^2}{(\conj{\zeta})^2}\conj{c_\zeta}
\conj{e^{2\ii\theta(\zeta)}}
=-\dfrac{q_0^2}{(\conj{\zeta})^2}\conj{C_\zeta}.
\end{align*}
完整四点留数逐一为
\begin{align*}
\Res_{z=\zeta}\dfrac{\mu_{+,2}(z)}{s_{22}(z)}
&=C_\zeta\mu_{-,1}(\zeta),\\
\Res_{z=\conj{\zeta}}\dfrac{\mu_{+,1}(z)}{s_{11}(z)}
&=-\conj{C_\zeta}\mu_{-,2}(\conj{\zeta}),\\
\Res_{z=h(\conj{\zeta})}\dfrac{\mu_{+,2}(z)}{s_{22}(z)}
&=-\dfrac{q_0^2}{(\conj{\zeta})^2}\conj{C_\zeta}
\mu_{-,1}(h(\conj{\zeta})),\\
\Res_{z=h(\zeta)}\dfrac{\mu_{+,1}(z)}{s_{11}(z)}
&=\dfrac{q_0^2}{\zeta^2}C_\zeta\mu_{-,2}(h(\zeta)).
\end{align*}
最后一行的系数来自
\[
-\conj{C_{\widehat\zeta}}
=-\conj{\left(-\dfrac{q_0^2}{(\conj{\zeta})^2}\conj{C_\zeta}\right)}
=\dfrac{q_0^2}{\zeta^2}C_\zeta.
\]
所以最后一个点的正号是两个负号相乘的结果.
全部新增展示式不另编号, 原文两条留数式及全部编号保持原样.
